\section{Results}
\label{sec:results}

\subsection{Proxy Similarity Is Necessary but Not Sufficient}
\label{sec:proxy_results}

On the four single-object Isaac Sim assets, the converted renders remain close
to the original MDL renders under standard proxy metrics: mean PSNR is
35.52\,dB, mean SSIM is 0.990, mean LPIPS is 0.020, CLIP cosine similarity is
0.925, and DINOv2 cosine similarity is 0.872 over 24 matched render pairs.
These measurements justify using the converted assets for deeper checks, but
they do not by themselves prove language-grounding preservation. The rest of
the paper therefore treats proxy similarity as a screening gate, not as the main
downstream-reliability claim. Figure~\ref{fig:render_scene_evidence_acl} shows the visual
counterpart to this gate and to the GRScenes grounding context, while the caption
keeps both visual panels below VLM-equivalence claims.

\begin{figure*}[t]
\centering
\includegraphics[width=\textwidth]{figures/fig_render_scene_evidence_wide.png}
\caption{Actual render evidence before task claims. Top: a representative
cropped single-object proxy detail pair. Bottom: a representative GRScenes
original/noMDL commercial-interior render pair. These panels orient readers, but
metric and task claims come from the frozen tables, not from the qualitative
figure.}
\label{fig:render_scene_evidence_acl}
\end{figure*}

\subsection{VLM Grounding Under a Frozen Material Perturbation}
\label{sec:vlm_grounding_pilot}

The clean GRScenes pool contains 15 visual-QA PASS original/noMDL pairs. It is
a pilot preservation control because it remains below the configured 20-pair
clean final gate. The lower row of
Figure~\ref{fig:render_scene_evidence_acl} anchors this section in actual
commercial-interior scene renders, while the numeric claims remain tied to the
frozen pair tables.

\begin{table*}[t]
\centering
\caption{Clean-pool GRScenes VLM grounding pilot over 15 blind-visual-QA PASS original/converted pairs. Counts are O/C hits under the structured-text prompt; the pool remains below the 20-pair benchmark threshold and is a control, not final benchmark evidence.}
\label{tab:grscenes_vlm_clean_pool_pass15}
\footnotesize
\setlength{\tabcolsep}{4pt}
\begin{tabular}{@{}>{\raggedright\arraybackslash}p{0.15\textwidth}
                >{\raggedright\arraybackslash}p{0.07\textwidth}
                >{\raggedright\arraybackslash}p{0.12\textwidth}
                >{\raggedright\arraybackslash}p{0.12\textwidth}
                >{\raggedright\arraybackslash}p{0.13\textwidth}
                >{\raggedright\arraybackslash}p{0.20\textwidth}@{}}
\toprule
\textbf{Probe} & \textbf{Rows} & \textbf{Answer} & \textbf{Raw point} & \textbf{Norm-1000 point} & \textbf{Scope} \\
\midrule
Gemma4 local & 30/30 & 15/15 / 15/15 & 0/15 / 0/15 & 8/15 / 6/15 & pilot preservation control \\
Qwen2.5-VL structured & 23/30 & 8/11 / 9/12 & 5/14 / 5/15 & 0/14 / 0/15 & coordinate-format diagnostic \\
\bottomrule
\end{tabular}
\par\vspace{0.25em}
\begin{minipage}{0.88\textwidth}
\scriptsize
\emph{Notes.} Gemma4 has 11/15 normalized-1000 same-hit-status pairs, 5/15
both-hit pairs, and 27.90 px mean point delta. Qwen2.5-VL normalized agreement
is an all-miss agreement over comparable points; raw pair agreement is 9/14
with 2/14 both-hit.
\end{minipage}
\end{table*}


Table~\ref{tab:grscenes_vlm_clean_pool_pass15} shows that answer-level
recognition can remain stable while point grounding is less stable. Gemma4
answers all 15 original and 15 converted renders correctly, but normalized-1000
point hits are 8/15 for original and 6/15 for converted. Qwen2.5-VL has lower
parser and coordinate stability: its structured-text run produces 23/30
scorable answer rows, raw point hits of 5/14 and 5/15, and zero normalized-1000
point hits in both material conditions.

The expanded stress set asks a different question. It freezes 30
target-centered PASS/WARN pairs where the target remains visible but material,
color, or lighting changes may be apparent.

\begin{table*}[t]
\centering
\caption{Expanded GRScenes material-shift grounding results over a frozen 30-pair target-centered stress set. Gemma4 is the canonical root run for this frozen set; Qwen2.5-VL is a second-model diagnostic under the same manifest. Values are hit counts for original/converted renders unless otherwise noted. Norm-1000 point-in-bbox is the primary point metric; raw point-in-bbox is retained only as a coordinate diagnostic.}
\label{tab:grscenes_vlm_stress_expanded30}
\scriptsize
\resizebox{\textwidth}{!}{%
\begin{tabular}{lllccccc}
\toprule
\textbf{Model} & \textbf{Role} & \textbf{Answer rows} & \textbf{Answer O/C} & \textbf{Raw point O/C} & \textbf{Norm-1000 point O/C} & \textbf{Norm hit agree} & \textbf{Norm both-hit} \\
\midrule
Gemma4 local & canonical & 60/60 & 30/30 / 30/30 & 1/30 / 1/30 & 27/30 / 29/30 & 28/30 & 27/30 \\
Qwen2.5-VL & diagnostic & 55/60 & 27/29 / 24/26 & 22/29 / 19/29 & 3/29 / 3/29 & 24/28 & 1/28 \\
\bottomrule
\end{tabular}%
}
\end{table*}


Table~\ref{tab:grscenes_vlm_stress_expanded30} reports the primary stress-set
result. Gemma4 preserves 30/30 target-category answers for both original and
converted renders. Under the primary normalized-1000 coordinate contract, it
scores 27/30 point hits on original renders and 29/30 on converted renders,
with 28/30 pairwise hit-status agreement. Qwen2.5-VL is different: it is more
comfortable in raw pixel space, with raw point hits 22/29 original and 19/29
converted, while normalized-1000 hits remain 3/29 for both conditions. This is
the main protocol finding: material perturbation, prompt contract, and
coordinate interpretation interact, so VLM grounding reliability cannot be
inferred from visual similarity alone.

\begin{table*}[t]
\centering
\caption{Reviewer-closure paired bootstrap confidence intervals for existing VLM and InternNav evidence. Deltas are converted/noMDL minus original MDL and are descriptive within the current frozen evidence pools, not population-level guarantees.}
\label{tab:reviewer_closure_paired_ci}
\scriptsize
\resizebox{\textwidth}{!}{%
\begin{tabular}{lllrrrr}
\toprule
\textbf{Evidence} & \textbf{Metric} & \textbf{N} & \textbf{Original} & \textbf{Converted} & \textbf{Delta [95\% CI]} & \textbf{Direction} \\
\midrule
GRScenes VLM expanded30 & Gemma4 Norm-1000 point hit & 30 & 0.9 & 0.9667 & 0.0667 [0.0, 0.1667] & converted higher \\
GRScenes VLM expanded30 & Gemma4 Raw point diagnostic & 30 & 0.0333 & 0.0333 & 0.0 [0.0, 0.0] & tie \\
GRScenes VLM expanded30 & Qwen2.5-VL Norm-1000 point hit & 28 & 0.1071 & 0.1071 & 0.0 [-0.1429, 0.1429] & tie \\
GRScenes VLM expanded30 & Qwen2.5-VL Raw point diagnostic & 28 & 0.7857 & 0.6429 & -0.1429 [-0.3214, 0.0] & converted lower \\
Official KuJiaLe InternNav val-unseen & Success rate & 99 & 0.5253 & 0.4848 & -0.0404 [-0.1515, 0.0707] & noMDL lower \\
Official KuJiaLe InternNav val-unseen & SPL & 99 & 0.4739 & 0.4298 & -0.0441 [-0.1417, 0.0514] & noMDL lower \\
Official KuJiaLe InternNav val-unseen & Navigation error & 99 & 3.6798 & 3.6306 & -0.0492 [-0.6464, 0.5386] & noMDL lower \\
Official KuJiaLe InternNav val-unseen & Trajectory length & 99 & 6.9754 & 7.0598 & 0.0844 [-0.7366, 0.9304] & noMDL higher \\
\bottomrule
\end{tabular}%
}
\end{table*}


The paired bootstrap intervals in Table~\ref{tab:reviewer_closure_paired_ci}
calibrate the same point. On the expanded 30-pair stress set, Gemma4 has a small
positive converted-minus-original normalized-point delta, while the Qwen raw
point diagnostic decreases on converted renders. These are descriptive
intervals over frozen evidence pools, not population-level guarantees.

We therefore keep the VLM-grounding claim metric-auditable: scene panels orient
readers, but answer and point-grounding claims come from frozen tables;
overlay-style qualitative panels need clean render logs and VLM-overlay
provenance.

\subsection{Material Effects and NVIDIA Baseline Boundaries}
\label{sec:material_effect_baseline}

We next audit whether material mechanisms explain conversion risk and whether
the NVIDIA Asset Converter baseline changes the interpretation.

% Auto-generated by build_material_effect_risk_matrix.py
\begin{table*}[t]
\centering
\caption{Material-effect claim-boundary matrix for the ConvertAsset versus NVIDIA baseline. The table separates effect bins with bounded selected qualitative evidence from supplemental failure or limitation cases; it is not a population-level error-rate estimate.}
\label{tab:material_effect_risk_matrix}
\scriptsize
\resizebox{\textwidth}{!}{%
\begin{tabular}{llll}
\toprule
Effect & Evidence status & Relative interpretation & Claim allowed \\
\midrule
clearcoat & supplemental visual/PXR fail & selected NVIDIA failure; ConvertAsset target retained & selected NVIDIA failure case \\
opacity / transparency & selected qualitative ready & bounded three-condition evidence & bounded static + selected qualitative \\
emission & selected qualitative ready & bounded three-condition evidence & bounded static + selected qualitative \\
procedural texture & supplemental visual/PXR fail & both conversions lose checker evidence & diagnostic limitation case \\
normal / bump & selected qualitative ready & bounded three-condition evidence & bounded static + selected qualitative \\
displacement / height & selected qualitative ready & bounded three-condition evidence & bounded static + selected qualitative \\
\bottomrule
\end{tabular}%
}
\end{table*}


Table~\ref{tab:material_effect_risk_matrix} is the claim-boundary artifact. For
opacity/transparency, emission, normal/bump, and displacement/height, the
selected GRScenes pool has bounded static evidence for
original MDL, ConvertAsset noMDL, and the NVIDIA PreviewSurface baseline. This
supports a covered-bin comparison, not a universal visual-quality win: static
gates do not certify that the MDL mechanism itself is preserved, so the
selected qualitative evidence in
Figure~\ref{fig:material_effect_baseline_qualitative_acl} remains bounded to
the checked cases. The
supplemental clearcoat case is a selected NVIDIA target-loss failure while
ConvertAsset retains an approximated PreviewSurface target. Procedural texture
is a limitation case for both converters because checker/base-color texture
preservation evidence is lost after conversion. The visual panel makes the
covered-bin comparison inspectable, but Table~\ref{tab:material_effect_risk_matrix}
remains the governing claim-boundary artifact.

\begin{figure*}[t]
\centering
\includegraphics[width=0.94\textwidth]{figures/fig_material_effect_baseline_qualitative.png}
\caption{Selected material-effect qualitative comparison for the four
GRScenes-covered bins. Columns compare original MDL, ConvertAsset noMDL, and
the NVIDIA PreviewSurface baseline after the original-MDL cells passed the
log-checked clean rerender/provenance gate. The panels support bounded visual inspection
for the selected covered bins in Table~\ref{tab:material_effect_risk_matrix};
they do not estimate a population-level NVIDIA failure rate or certify
mechanism-level MDL preservation.}
\label{fig:material_effect_baseline_qualitative_acl}
\end{figure*}

\subsection{Embodied-Data Sanity and Official Scene Stability}
\label{sec:embodied_sanity}

For downstream embodied-data sanity, the official KuJiaLe
\texttt{val\_unseen} route covers 99 paired InternNav/DualVLN episodes across
three local official scenes. Combined original/noMDL means are SR
0.5253/0.4848, SPL 0.4739/0.4298, NE 3.6798/3.6306, and TL 6.9754/7.0598. The
paired confidence intervals in Table~\ref{tab:reviewer_closure_paired_ci}
cross zero for SR, SPL, NE, and TL, so the supported conclusion is scoped:
converted scenes can be evaluated end-to-end in a real official navigation
stack, but the current evidence does not establish broad embodied-navigation
robustness. Selected rollout stills and videos are retained as supplemental
qualitative evidence for figure and rebuttal use; they do not replace the
99-episode paired run.

\begin{table*}[t]
\centering
\footnotesize
\setlength{\tabcolsep}{2pt}
\begin{tabular}{@{}>{\raggedright\arraybackslash}p{0.30\textwidth}>{\raggedright\arraybackslash}p{0.30\textwidth}>{\raggedright\arraybackslash}p{0.30\textwidth}@{}}
\toprule
Metric & Original MDL & noMDL \\
\midrule
Scenes & 3 & 3 \\
Successful runs & 9; 0 fail & 9; 0 fail \\
Ready time (s) & 13.95 [11.64, 16.28] & 14.12 [12.31, 16.24] \\
FPS & 2.36 [2.01, 2.70] & 2.42 [2.08, 2.73] \\
\bottomrule
\end{tabular}
\caption{Official KuJiaLe / InteriorAgent scene load/render stability. Aggregate original/noMDL rows only; intervals are bootstrap 95\% CIs. The NVIDIA official-scene row is omitted until matching converted scenes pass the same load/render check.}
\label{tab:official_scene_performance_summary}
\end{table*}


Table~\ref{tab:official_scene_performance_summary} adds the official-scene
load/render stability check. All 18 required paired runs complete successfully
in fresh processes across original and noMDL conditions on the same three
official scenes, i.e., 18/18 successful required runs. Aggregated
warmed-ready time is comparable between original MDL and noMDL, with
overlapping confidence intervals. We therefore use this evidence for
loadability and stability only. We do not report an NVIDIA official-scene
performance row because matching official-scene NVIDIA conversions have not yet
been generated and smoke-gated.
